\section{Introduction}
\label{sec:introduction}

Modular robots promise a form of embodied adaptability that a conventional
fixed-morphology robot cannot provide. A collection of repeated modules can change
topology, redistribute actuation, form task-specific kinematic chains, recover from
failures, or separate into independently mobile agents. The configuration of the
robot is therefore not merely a parameter of its controller; it is a time-varying
part of the system state. This observation has shaped the field of modular
self-reconfigurable robotics for decades \citep{yim2007modular}.

Simulation infrastructure has not always followed the same abstraction. General
robotics engines accept bodies, joints, actuators, collision geometry, and constraints.
Those primitives are necessary, but they do not answer modular-robot questions such
as:

\begin{itemize}
  \item Which links and joints constitute one reusable module type?
  \item Which physical surfaces are connectors, and which pairs may mate?
  \item What translational, angular, roll, and speed errors are acceptable at capture?
  \item When is an engine constraint considered a committed logical connection?
  \item Which modules form an assembly after a dock or release?
  \item How should a changing morphology be exposed to planners, algorithms, logs, and
        graphical tools?
\end{itemize}

In a simulator-specific prototype, these concepts are often distributed across model
files, Python scripts, controller constants, scene-construction code, and assumptions
inside a viewer. That approach can produce an effective demonstration, but it creates
three long-term problems. First, semantic knowledge is not portable: replacing one
engine with another requires reconstructing the meaning of connectors and
configurations. Second, physical and logical state can diverge: an application may add
an edge to its topology graph before the engine has actually created a constraint, or
may remove a graph edge while a stale constraint remains active. Third, the tooling
surface becomes fragile: authoring, validation, simulation, algorithms, and
visualization each invent a slightly different representation of the same robot.

\modsim{} addresses this gap with a semantics-first architecture. It is not a new
rigid-body solver. Instead, it defines a portable \robotpack{}, a canonical runtime,
and narrow interfaces to physics engines and user interfaces. The central design rule
is:

\begin{quote}
\emph{ModSim owns modular-robot meaning and canonical history; a backend owns physics;
every visualization or algorithmic view is derived from the canonical state.}
\end{quote}

This separation makes the framework useful even before multiple high-fidelity backends
exist. A dependency-free mock adapter can exercise connector semantics and event
ordering in continuous integration. A \mujoco{} adapter can then supply articulated
dynamics, contact, effort actuation, runtime welds, and native visualization without
moving connector rules into engine-specific code. Future adapters---for example an
Isaac Sim implementation---can be checked against the same semantic contract.

\subsection{Research questions}

The current implementation is organized around four questions.

\begin{enumerate}
  \item \textbf{Semantic portability.} Can one platform description define modules,
        connectors, docking policy, metadata, and generated model views without
        depending on one simulator's scene format?
  \item \textbf{Runtime consistency.} Can physical constraints, logical connections,
        assembly membership, events, and graph views remain coherent during docking
        and undocking?
  \item \textbf{Progressive physical fidelity.} Can a generic semantic runtime begin
        with kinematic tests and then support real gravity, contact, articulated joint
        control, and multi-module locomotion without replacing its state model?
  \item \textbf{Inspectability.} Can authors and runtime users validate the same data,
        observe the same event history, and view physical and topological state from a
        single authoritative simulation?
\end{enumerate}

\subsection{Contributions}

This technical report describes four principal contributions.

\begin{enumerate}
  \item A strict, self-contained \robotpack{} format that separates mechanical assets
        from modular semantics, backend mappings, custom metadata, capabilities, and
        model-view recipes. The associated loader, validator, canonical writer, URDF
        importer, and Studio application provide an end-to-end authoring workflow.
  \item A backend-independent runtime that resolves connector frames, evaluates
        compatibility and geometric acceptance, applies policy guards, performs a
        two-phase physical/logical commit, evolves assemblies, accepts atomic joint
        command batches, and records deterministic semantic events and metrics.
  \item Narrow physics and visualization boundaries: a reference mock backend, an
        optional \mujoco{} backend, immutable model-view data-transfer objects, a
        revision-aware view factory, and a dual-window Runtime Inspector driven by one
        authoritative session.
  \item A \smoresep{} case study containing committed visual and collision assets,
        four docking faces, a physics-specific MJCF model, wheel and joint effort
        control, tire contact, a physical two-module docking cycle, and a seven-module
        wheel-driven realization of a published Driver-to-Snake topology plan.
\end{enumerate}

\subsection{Scope and non-claims}

The implementation is pre-alpha software (version 0.1.0), and its strongest evidence
is software and simulation evidence. The report does \emph{not} claim a new contact
solver, autonomous reconfiguration planner, calibrated \smoresep{} digital twin,
magnetic-field model, or reproduction of a published Driver-to-Snake hardware
trajectory. The seven-module routes, friction parameters, effort limits, skid model,
and control gains are explicit simulation choices. The published 2019 work supplies
the topology and connector substitutions; \modsim{} supplies one physically plausible,
deterministic execution of that plan.

\subsection{Organization}

\Cref{sec:background} reviews relevant modular-robot and simulation concepts.
\Cref{sec:requirements} turns the observed gap into requirements and authority rules.
\Cref{sec:architecture} presents the overall architecture. The following four sections
describe Robot Packs and Studio, runtime docking, backend physics, and model views with
runtime visualization. \Cref{sec:smores} develops the \smoresep{} case study and its
controllers. \Cref{sec:experiments,sec:analysis} report experiments and analyze the
evidence. \Cref{sec:future} states limitations and a staged roadmap. The appendices
provide schema summaries, event taxonomies, scenario commands, and a reproducibility
checklist.

